\begin{table}[tp]\centering\footnotesize\setlength{\tabcolsep}{3pt}
\caption{Implied Trade-Off: The Phillips Multiplier}\label{tab:tradeoff}
\resizebox{\ifdim\width>\textwidth\textwidth\else\width\fi}{!}{%
\begin{tabular}{lcccc}\toprule
 & \multicolumn{2}{c}{24 months} & \multicolumn{2}{c}{36 months}\\ \cmidrule(lr){2-3}\cmidrule(lr){4-5}
Rates and FCI, measured as & Estimate & 90\% set & Estimate & 90\% set\\ \midrule
Nominal, 1996 start & $-$1.02 & [$-$3.74, $-$0.09] & $-$0.76 & [$-$2.46, 0.03]\\
\textbf{Real, 1996 start (baseline)} & \textbf{$-$0.85} & \textbf{[$-$1.61, $-$0.26]} & \textbf{$-$0.61} & \textbf{[$-$1.23, $-$0.11]}\\
Nominal, 1973 start & $-$1.11 & undefined; 68\%: [$-$2.48, $-$0.08] & $-$1.05 & undefined; 68\%: [$-$2.96, 0.00]\\
Real, 1973 start & $-$1.85 & [$-$5.36, $-$0.77] & $-$1.52 & [$-$5.82, $-$0.49]\\
\bottomrule\end{tabular}}
\begin{flushleft}\footnotesize\textit{Notes:} $\kappa_H$ is the core PCE price level at month $H$ (percent) per percentage-point year of cumulative unemployment-gap response \citep{barnichon2021}; its inverse in absolute value is the sacrifice ratio. Estimate: ratio of the point responses. 90\% set: Anderson--Rubin type, all values for which the 90\% bootstrap interval of $\tilde P_H-\kappa\tilde U_H$ contains zero. The ratio is undefined when the 90\% interval of the cumulative unemployment response includes zero, because the data then cannot rule out unemployment responses of zero or either sign; the 68\% set is then shown where it exists. 12 lags. All cases in Appendix Table~\ref{tab:tradefull}.\end{flushleft}
\end{table}
